\section{Relation to earlier work}\label{sec:related-work}

The cubic cover used here belongs to an established family.
Kazalicki--Scholl \cite[Sections~7--8]{KS} describe the genus-zero
groups $\Phi_0(N)$ and the elliptic family
$Y^2=X(X+1)(X+t^N)$. For $N=3$ and $t=z$, substituting $x=-X$
in our equation gives
\[
 y^2=-X(X+1)(X+z^3).
\]
Thus the two families differ by the constant quadratic twist $-1$.
The corresponding $\operatorname{Sym}^4$ coefficient systems are
isomorphic over $\mathbb Q$, since the fourth power of the twisting
character is trivial. Their analyzed weight-three representation uses
$R^1$; it does not by this comparison identify the fourth-symmetric-power constituent
studied here. Neither the underlying family nor its cubic-cover geometry
is claimed as new.

Scholl's construction places parabolic cohomology of symmetric powers
in the general theory of motives attached to noncongruence cusp forms;
see \cite[Introduction]{Scholl}. More precise projector references are
\cite[Sections~1.1--1.2]{SchollMotives} and
\cite[Sections~2.2--2.8 and 4.2--4.4]{SchollVanishing}.
Our local calculations use the standard
etale-cohomological framework. Section~\ref{sec:scholl} applies these
established constructions with an explicit full-level cover and projectors
matching the literal normalized deck and inversion operators. The degree-four
representations treated by Atkin--Li--Liu--Long \cite{ALLL} provide
further relevant context. Their automorphy theorem requires quaternion
multiplication satisfying specified compatibility conditions. We have
not established those conditions for $W_4^+$ or $W_4^-$
\cite[Definition~3.1.1 and Theorem~4.3.1]{ALLL}. In fact
$J=A-A^2$ exchanges the two spaces, while $I$ is scalar on each;
these operators are not a quaternion-multiplication pair on either
four-dimensional summand.

The auxiliary moment identities also have substantial precedents.
The quadratic-twist corollaries use the Legendre Hecke trace formulas
recorded in \cite[Theorem~2.3]{Grove}, with their attribution to Ahlgren
and Ahlgren--Ono \cite{Ahlgren,AhlgrenOno}. The weight-four forms used in the
second-moment comparisons occur in Long--Tu--Yui--Zudilin
\cite[Theorem~2 and Table~1]{LTYZ}; those modular forms and the
associated hypergeometric modularity are not discoveries of this note.

There is also an exact parameter overlap with Allen--Grove--Long--Tu
\cite{EHMM2}. The datum in our cubic second-moment normalization,
\[
 \left\{\begin{matrix}
 \tfrac12,\tfrac12,\tfrac13,\tfrac13\\
 1,1,\tfrac56,\tfrac56
 \end{matrix}\right\},
\]
is their $\mathrm{HD}_4(j/12)$ at $j=4$, defined in the published
article's introduction immediately before equation~(1.1).
(The corresponding definition is equation~(1.2) in arXiv v1.)
Their Theorem~1.4 concerns its specialization at argument $-1$;
our normalization identity uses argument $1$ and an explicit squared
Jacobi-sum factor. Equality of the parameter lists alone does not
identify these specializations or identify either with $W_4^\pm$.

The results established here specify the arithmetic coefficient sheaf,
construct its two rational descents, and derive their good-prime traces,
including the inert-prime boundary sign. These statements concern the
precisely normalized objects defined in the paper. Their exact priority
has not been established by the literature comparison, and no priority
claim is made for the constituent or trace formula. The conductor calculation in Sections~\ref{sec:local-three}--\ref{sec:conductors}
is an independent local argument, not a consequence of this comparison.
No automorphic identification or classical elliptic BSD consequence
follows from the literature overlap.
